\documentclass[12pt]{article}
\usepackage{amsmath,amssymb,amsfonts}
\usepackage[margin=1in]{geometry}
\usepackage{bm}

\newcommand{\vect}[1]{\mathbf{#1}}
\newcommand{\grad}{\boldsymbol{\nabla}}
\newcommand{\curl}{\grad\times}
\newcommand{\divr}{\grad\cdot}

\begin{document}

\section*{Problem 9}

\textbf{Problem:} A physical system known as a \emph{vector meson} is described by the three vector fields $\vect{E}$, $\vect{H}$, $\vect{A}$ and a scalar field $V$, which satisfy the following equations:
\begin{align}
\curl\vect{E} &= -\mu^2\vect{A} \qquad (\mu^2 \text{ is a positive constant}), \label{eq:curlE} \\
\vect{H} &= \curl\vect{A}, \label{eq:H} \\
\vect{E} &= -\grad V, \label{eq:E} \\
\divr\vect{E} &= \mu^2 V. \label{eq:divE}
\end{align}

\begin{itemize}
\item[(a)] Show that $\divr\vect{H} = 0$.
\item[(b)] Show that $V$ satisfies the equation: $\nabla^2 V - \mu^2 V = 0$.
\item[(c)] When $V$ is time-independent, it satisfies $\nabla^2 V - \mu^2 V = 0$. If this is the case inside a closed surface and if $V$ vanishes on the closed surface, show that $V = 0$ everywhere inside.
\end{itemize}

\bigskip
\textbf{Solution:}

\subsection*{(a) $\divr\vect{H} = 0$}

From Eq.~\eqref{eq:H}:
\[
\divr\vect{H} = \divr(\curl\vect{A}).
\]

By the vector identity (Eq.~1.84 in Byron \& Fuller): \emph{the divergence of any curl is identically zero}:
\[
\divr(\curl\vect{A}) \equiv 0.
\]

Therefore:
\[
\boxed{\divr\vect{H} = 0}. \quad \blacksquare
\]

\subsection*{(b) $\nabla^2 V - \mu^2 V = 0$}

From Eq.~\eqref{eq:E}, $\vect{E} = -\grad V$. Taking the divergence:
\[
\divr\vect{E} = -\nabla^2 V.
\]

But from Eq.~\eqref{eq:divE}, $\divr\vect{E} = \mu^2 V$. Equating:
\[
-\nabla^2 V = \mu^2 V,
\]
\[
\boxed{\nabla^2 V - \mu^2 V = 0}. \quad \blacksquare
\]

\textit{Remark:} This is the \emph{time-independent Klein--Gordon equation} (or Yukawa equation). Its spherically symmetric solution is $V \propto e^{-\mu r}/r$, the Yukawa potential describing the short-range nuclear force mediated by mesons.

\subsection*{(c) If $V$ satisfies $\nabla^2 V - \mu^2 V = 0$ inside a closed surface $S$, and $V = 0$ on $S$, then $V = 0$ everywhere inside}

Multiply the equation $\nabla^2 V - \mu^2 V = 0$ by $V$ and integrate over the volume $\tau$ enclosed by $S$:
\[
\int_\tau V\,\nabla^2 V\;d\tau - \mu^2\int_\tau V^2\;d\tau = 0.
\]

Apply Green's first identity (from Problem 7) with $\varphi = \psi = V$:
\[
\int_\tau V\,\nabla^2 V\;d\tau = \oint_S V\,\grad V\cdot d\boldsymbol{\sigma} - \int_\tau |\grad V|^2\;d\tau.
\]

Since $V = 0$ on $S$, the surface integral vanishes. Substituting:
\[
-\int_\tau |\grad V|^2\;d\tau - \mu^2\int_\tau V^2\;d\tau = 0.
\]

Therefore:
\[
\int_\tau \bigl[|\grad V|^2 + \mu^2 V^2\bigr]\;d\tau = 0.
\]

The integrand is a sum of two non-negative terms (since $\mu^2 > 0$). A non-negative function integrating to zero must vanish everywhere:
\[
|\grad V|^2 = 0 \quad \text{and} \quad V^2 = 0 \quad \text{throughout } \tau.
\]

Therefore:
\[
\boxed{V = 0 \text{ everywhere inside } S}. \quad \blacksquare
\]

\textit{Note:} This proof relies crucially on $\mu^2 > 0$. For Laplace's equation ($\mu^2 = 0$), we can only conclude $|\grad V|^2 = 0$, hence $V = \text{const}$, and the boundary condition $V|_S = 0$ then gives $V = 0$. But for the Helmholtz equation $\nabla^2 V + k^2 V = 0$ (with the opposite sign), the argument fails---nontrivial solutions with $V = 0$ on the boundary do exist (these are the eigenmodes).

\end{document}
